% Auto-generated by he-bench-paper/tables/T06_model_effects/render.py -- do not edit by hand.
\begin{table}[t]
\centering\small
\setlength{\tabcolsep}{4pt}
\begin{tabular}{@{}l rc@{}}
\toprule
& \multicolumn{2}{c}{\textbf{LSS-$\lambda$}}\\
\cmidrule(l){2-3}
\textbf{Model} & gain units & tier\\
\emph{resolution} & $\pm0.058$ & \\
\midrule
claude-opus-5 & $+0.228$ & 1\\
claude-sonnet-5 & $+0.029$ & 2\\
kimi-k3 & $-0.014$ & 2\\
gpt-5.6-sol & $-0.069$ & 2\\
gpt-5.6-terra & $-0.174$ & 3\\
\bottomrule
\end{tabular}
\caption{\textbf{The model effect on the balanced scope.} LSS-$\lambda$ is the model term of an additive fit over task and model, in normalized-gain units: what swapping the optimizer model is worth once the task is accounted for. This is the numeric form of the right panel of Figure~\ref{fig:showdown}, from the same call, so the two cannot disagree. \emph{resolution} is the estimator's own measured split-round swing and the tier column cuts at it, so a shared tier means ``closer together than re-running the grid moves them'' and ranks within a tier are not resolved. Two further estimators of the same quantity --- gain standardized within task, and mean within-group rank --- are computed but not shown; all three produce the identical ordering (pairwise rank concordance $\geq1.00$), so the ranking here does not depend on the additive assumption. Scope: 15 contestant rows over 3 tasks on \oc\ alone, balanced. GAIA is excluded because its measured-zero baseline makes its gain a raw held-out score rather than a fraction of headroom. No harness effect is reported: this scope holds one harness, so the fit has no harness factor to estimate --- the harness magnitude in Section~\ref{sec:rq1} comes from the paired contrast instead.}
\label{tab:model-effects}
\end{table}
